\section{Evaluation methodology}
\label{sec:methodology}

\subsection{Commitments}

\begin{itemize}
  \item \textbf{Cross-dataset evaluation.} Domain shift is the most likely
        cause of real-world failure; a random split of the training set
        flatters a model and predicts nothing about deployment.
  \item \textbf{Report the gap.} In-domain and out-of-domain figures appear
        side by side.
  \item \textbf{Event-level evaluation on sequences}, not frame-level on
        stills. The validator only demonstrates value over time.
  \item \textbf{$F_2$ rather than $F_1$ as the headline}, for the reason given
        in \cref{eq:fbeta-ratio}.
  \item \textbf{Never quote a figure we did not measure.} Published
        third-party numbers appear only in clearly attributed rows.
  \item \textbf{Every proportion carries its interval} (\cref{sec:intervals}).
\end{itemize}

\begin{recorded}{A trap we walked into and documented}
Our first false-alarm measurement on ignition sequences gave 0.8\,\%, against
71\,\% on curated hard negatives with the same model and threshold. The
difference is negative-set difficulty: frames shortly before a fire are
usually clear skies, whereas curated negatives are cloud, fog and dust chosen
to confuse. Had we published the 0.8\,\% figure as evidence the validator
works, any reviewer inspecting the negative set would have dismantled the
claim. Every false-alarm figure in this document names the negative set it
came from.
\end{recorded}

\subsection{Reproducibility}

Detection is run once and cached to JSONL; validator configurations are then
replayed against the cache. Every configuration therefore sees byte-identical
detector output, so any difference in outcome is attributable to the validator
and not to run-to-run variation. All results files are committed under
\texttt{eval/results/}, and every reported figure traces to one of them
(\cref{sec:reproducibility}).

A third guard, \texttt{tools/freeze\_results.py}, checks each published figure
against the results file it cites and fails the build if the two disagree.

\begin{recorded}{The freeze guard catching its author}
The aerial flood model was first registered with its figures marked
\texttt{measured}, the same label the held-out evaluations carry, with no
source file named. The guard rejected it. The entry now reads
\texttt{validation\_measured} with an explicit note that FloodNet's withheld
masks make an independent split impossible — which is the distinction
\cref{sec:flood} now draws in the text as well.
\end{recorded}

\subsection{Execution backends}
\label{sec:backends}

Two backends exist and are not interchangeable. \textsc{reference} runs the
as-published dynamic-shape graph on the CPU provider and is deterministic;
every figure in this document comes from it. \textsc{fast} runs a static-shape
variant on Apple's CoreML provider and is 1.7--1.9$\times$ faster
(\SI{34}{\milli\second} against \SI{60}{\milli\second} at \SI{1024}{\pixel}).

\textsc{fast} is not bit-identical. Across 100 validation frames the detection
count matched on 100 of 100 and boxes agreed to within \SI{0.41}{\pixel}, but
confidences drift by up to $0.012$ because the Neural Engine computes at lower
precision. At the $0.20$ operating point that drift changed no metric; at
$0.10$ and $0.30$ it moved box precision and recall by about one percentage
point by flipping a single detection across the threshold. \textsc{fast} is
therefore used for anything a person watches, and \textsc{reference} for
anything reported.

\begin{recorded}{A parity check that checked nothing}
The ONNX export parity check originally traced the model on random noise. A
trained detector finds nothing in random noise, so the check compared zero
detections against zero detections and passed unconditionally. It now traces a
real frame and reports \textsc{unverified} if that frame yields no detections,
instead of passing quietly.
\end{recorded}

\subsection{Demonstration material}
\label{sec:demo}

Demonstration clips are treated as evidence, not decoration, and are
subject to the same gating as any published figure. One clip per hazard is
assembled from local data in sorted order at a fixed frame rate, so a clip
built twice is identical and a run over it reproduces exactly. The recording
is rendered directly from the pipeline instead of screen-captured, which
removes any dependence on a browser, a display or a network at presentation
time.

Two independent conditions decide whether a clip may appear in a recorded
demonstration, a figure or any public material. Both are recorded per clip in
a manifest, and the renderer refuses a clip that fails either — refusing is
the safe default, because a clip that should not be shown plays exactly like
one that may.

\begin{enumerate}
  \item \textbf{The licence permits publication.} The source datasets carry
        four different licences and one forbids redistribution outright. Where
        a licence permits publication subject to attribution, the credit is
        burned onto every frame not left to a caption that may be
        cropped away.
  \item \textbf{The frames are a genuine time sequence.} The validator is
        temporal, so demonstrating it on footage without temporal continuity
        misrepresents it in both directions: persistence rejects almost
        everything because nothing genuinely persists, while a few detections
        confirm anyway because unrelated objects happen to overlap between
        consecutive frames and the tracker associates them.
\end{enumerate}

\begin{recorded}{A clip that played perfectly and meant nothing}
The first people-in-water clip was 60 unrelated validation stills stitched
into a video. It played correctly and produced confident figures — 156
detections proposed, 16 confirmed, 121 rejected by persistence — all of them
artefacts. Continuity is now verified, not assumed: adjacent-frame
histogram correlation is about $0.997$ on genuine footage against about
$0.015$ for unrelated frames from the same split, and the fetch script refuses
to report success unless that separation holds. On real sequential footage the
same detector and validator give 451 proposed and 59 confirmed.
\end{recorded}

\begin{recorded}{A sound method on an unrepresentative sample}
DRespNeT is published as shuffled, renamed stills whose filenames retain the
original capture index, and ordering by that index does recover the authors'
order. A twelve-frame window looked continuous — adjacent correlation $0.695$
against $0.235$ for unrelated pairs — and a 66-frame clip was built on that
basis and described as a genuine sequence. That was wrong. The authors also
publish their un-augmented raw frames, and across all 615 of them the median
adjacent correlation is $0.388$ with 362 scene cuts; the longest continuous
run in the entire dataset is \emph{seven} frames, and within the 66-frame
range actually used, 40 of 65 transitions are scene cuts. The check was the
right kind of check — asking whether detections associate across frames rather
than whether frames look alike — which is what made the conclusion feel
well-founded. Tracks did reach length seven and 124 associations survived to
persistence three, but on footage that is 62\,\% scene cuts those associations
are coincidental spatial overlap between unrelated scenes. The claim was
retracted; the clip is retained for exercising the detector only. Sample size
is the part that has to be checked first.
\end{recorded}

\begin{table}[htbp]
\centering\small
\caption{Demonstration footage and its licence. Every clip in the
demonstration path is publicly licensed; no restricted footage remains.}
\label{tab:clips}
\begin{tabularx}{\textwidth}{l L l}
\toprule
Hazard & Source & Licence \\
\midrule
Fire            & HPWREN FIgLib ignition sequence~\citep{hpwren} & attribution required \\
Flood           & USGS gauge camera, real flood event~\citep{usgshivis} & public domain \\
People in water & SeaDronesSee flight sequence~\citep{varga2022seadronessee} & CC0 1.0 \\
Building access & DRespNeT stills — \emph{no sequence exists}~\citep{sirma2025drespnet} & CC BY 4.0 \\
\bottomrule
\end{tabularx}
\end{table}

Flood was solved by changing source, not by relaxing the rule. The USGS
operates over 1{,}300 streamgage cameras whose imagery is unrestricted, and
they are the right kind of footage as well as the right licence: fixed mount,
pointed at water, sampled roughly hourly — the same sparse regime the trend
detection was calibrated for, and the regime in which a fixed
\SI{300}{\second} window originally failed outright (\cref{sec:flood-trend}).
Only daylight frames are used: these cameras switch to infrared at night, and
on night frames the segmenter tints sky and vegetation as water. That is the
detector's declared envelope, and the fetch script enforces it instead of
restating it.
